\documentclass[12pt,a4paper]{article}

\usepackage[utf8]{inputenc}
\usepackage[T1]{fontenc}
\usepackage[romanian]{babel}
\usepackage{geometry}
\usepackage{setspace}
\usepackage{array}
\usepackage{booktabs}
\usepackage{longtable}
\usepackage{hyperref}

\geometry{margin=2.5cm}
\onehalfspacing

\title{\textbf{Analiza Comparativă a Algoritmilor de Sortare}}
\author{
    Proiect: Scriere Academica \\
    Universitatea de Vest din Timișoara \\
    Facultatea de Informatică \\
    Bosioc Aser \\
    \texttt{aser.bosioc07@e-uvt.ro}
}
\date{2026}

\begin{document}

\maketitle
\thispagestyle{empty}
\newpage

\section{Introducere}

Această lucrare prezintă o analiză comparativă a mai multor algoritmi de sortare implementați în limbajul C++. Scopul principal al proiectului este de a evidenția diferențele dintre metodele de sortare din punct de vedere al timpului de execuție și al memoriei auxiliare utilizate. În practică, alegerea algoritmului potrivit poate influența semnificativ performanța unei aplicații, mai ales atunci când volumul de date este mare sau când datele au o anumită structură.

Sortarea este o operație fundamentală în informatică. Ea apare în numeroase contexte: organizarea datelor, optimizarea căutărilor, preprocesarea informației pentru algoritmi mai complecși, generarea de rapoarte sau prelucrarea colecțiilor mari de valori. Din acest motiv, analiza algoritmilor de sortare reprezintă un subiect important atât în formarea teoretică, cât și în activitatea practică de programare.

În cadrul acestui proiect au fost implementați nouă algoritmi de sortare: Bubble Sort, Insertion Sort, Selection Sort, Merge Sort, Quick Sort, Heap Sort, Shell Sort, Radix Sort și IntroSort. Pentru fiecare algoritm au fost măsurate timpul mediu de execuție și memoria estimată, folosind mai multe tipuri de date de intrare și mai multe dimensiuni de vectori.

\section{Scopul proiectului}

Prin această lucrare s-au urmărit mai multe obiective:

\begin{itemize}
    \item implementarea într-un singur proiect a unor algoritmi de sortare clasici și moderni;
    \item compararea lor în condiții de test identice;
    \item observarea influenței tipului de date asupra performanței;
    \item analiza consumului de memorie pentru fiecare metodă;
    \item obținerea unei documentații academice care să explice atât teoria, cât și rezultatele experimentale.
\end{itemize}

\section{Descrierea aplicației}

Aplicația este implementată în fișierul \texttt{sort/main.cpp} și are rolul de a genera date, de a executa benchmark-ul și de a salva rezultatele în fișierul \texttt{sort/results.csv}. Programul folosește funcționalități din biblioteca standard C++, în special \texttt{std::chrono}, pentru măsurarea timpilor de execuție.

Pentru fiecare combinație dintre algoritm, tip de date și dimensiune a intrării, programul:

\begin{enumerate}
    \item generează vectorul inițial;
    \item copiază datele într-un vector de lucru;
    \item rulează algoritmul de sortare de mai multe ori;
    \item calculează timpul mediu de execuție;
    \item verifică dacă rezultatul final este corect sortat;
    \item estimează memoria suplimentară utilizată.
\end{enumerate}

Această abordare permite obținerea unor rezultate coerente și comparabile, reducând influența erorilor de măsurare la dimensiuni mici.

\section{Algoritmii implementați}

\subsection{Bubble Sort}

Bubble Sort este unul dintre cei mai simpli algoritmi de sortare. El compară elementele vecine și le interschimbă dacă acestea nu sunt în ordinea corectă. Procesul se repetă până când nu mai este necesară nicio interschimbare.

Avantajul principal este simplitatea implementării, însă performanța este slabă pentru date mari. Complexitatea medie și nefavorabilă este de ordinul $O(n^2)$, iar memoria suplimentară este $O(1)$.

\subsection{Insertion Sort}

Insertion Sort construiește progresiv o porțiune sortată a vectorului. Fiecare element nou este inserat în poziția corespunzătoare printre elementele deja ordonate.

Acest algoritm este eficient pentru vectori mici sau pentru date aproape sortate. În cel mai bun caz are complexitate $O(n)$, dar în medie și în caz nefavorabil are complexitate $O(n^2)$. Memoria suplimentară este constantă.

\subsection{Selection Sort}

Selection Sort selectează la fiecare pas minimul din partea nesortată a vectorului și îl mută pe poziția corectă. Algoritmul este ușor de înțeles, dar nu oferă performanțe bune în practică pentru volume mari de date.

Complexitatea rămâne $O(n^2)$ aproape indiferent de distribuția datelor, iar memoria suplimentară este $O(1)$.

\subsection{Merge Sort}

Merge Sort aplică strategia divide et impera. Vectorul este împărțit în subvectori mai mici, care sunt sortați recursiv și apoi interclasați.

Acest algoritm este predictibil și stabil, având complexitate $O(n \log n)$ în toate cazurile importante. Dezavantajul său principal constă în memoria suplimentară de ordin $O(n)$, necesară pentru interclasare.

\subsection{Quick Sort}

Quick Sort alege un pivot și împarte vectorul în două părți: elemente mai mici decât pivotul și elemente mai mari decât pivotul. Procedura este aplicată recursiv pe subvectorii rezultați.

În practică, Quick Sort este foarte rapid, având performanțe excelente pe multe tipuri de date. Complexitatea medie este $O(n \log n)$, însă în unele cazuri nefavorabile poate ajunge la $O(n^2)$. Memoria suplimentară este aproximativ $O(\log n)$.

\subsection{Heap Sort}

Heap Sort se bazează pe structura de date heap. Mai întâi transformă vectorul într-un heap, apoi extrage repetat maximul și îl mută la finalul vectorului.

Avantajul său este că garantează complexitate $O(n \log n)$ fără memorie auxiliară mare. În schimb, în practică poate fi mai lent decât Quick Sort sau IntroSort.

\subsection{Shell Sort}

Shell Sort este o extindere a lui Insertion Sort. În loc să compare doar elemente vecine, compară elemente aflate la anumite distanțe, numite \emph{gap}-uri, iar aceste distanțe sunt reduse treptat până la $1$.

În acest proiect, Shell Sort reprezintă una dintre cele trei metode noi adăugate. El oferă performanțe mai bune decât algoritmii clasici pătratici și se comportă bine pentru dimensiuni medii de intrare. Memoria suplimentară rămâne $O(1)$.

\subsection{Radix Sort}

Radix Sort este un algoritm necomparativ, potrivit pentru numere întregi. Implementarea folosește sortarea pe octeți și tratează corect și numerele semnate.

Avantajul său principal constă în viteza foarte bună pentru seturi mari de valori numerice. Complexitatea este de forma $O(k \cdot n)$, unde $k$ reprezintă numărul de pași necesari pentru parcurgerea cifrelor sau a octeților. Dezavantajul este consumul suplimentar de memorie, de ordin $O(n)$.

\subsection{IntroSort}

IntroSort este un algoritm hibrid, reprezentat în proiect prin \texttt{std::sort}. Acesta combină avantajele mai multor metode și evită degradările severe de performanță caracteristice unor implementări simple de Quick Sort.

IntroSort este foarte important în proiect, deoarece reflectă o soluție practică, folosită frecvent în aplicațiile reale. Complexitatea sa este $O(n \log n)$, iar memoria suplimentară este aproximativ $O(\log n)$.

\section{Tabel comparativ}

\begin{longtable}{>{\raggedright\arraybackslash}p{3cm} >{\raggedright\arraybackslash}p{3cm} >{\raggedright\arraybackslash}p{3cm} >{\raggedright\arraybackslash}p{4cm}}
\toprule
Algoritm & Tip & Complexitate medie & Observație principală \\
\midrule
Bubble Sort & comparații & $O(n^2)$ & foarte lent pentru date mari \\
Insertion Sort & comparații & $O(n^2)$ & bun pentru date aproape sortate \\
Selection Sort & comparații & $O(n^2)$ & simplu, dar ineficient \\
Merge Sort & divide et impera & $O(n \log n)$ & stabil și predictibil \\
Quick Sort & divide et impera & $O(n \log n)$ & foarte rapid în practică \\
Heap Sort & heap & $O(n \log n)$ & robust și cu memorie mică \\
Shell Sort & inserare cu gap & sub-pătratic în practică & bun pentru dimensiuni medii \\
Radix Sort & necomparativ & $O(k \cdot n)$ & excelent pentru întregi \\
IntroSort & hibrid & $O(n \log n)$ & foarte bun pentru uz general \\
\bottomrule
\end{longtable}

\section{Metodologia experimentală}

Pentru a obține o imagine cât mai clară asupra comportamentului algoritmilor, testele au fost realizate pe mai multe tipuri de date:

\begin{itemize}
    \item \texttt{Random} -- valori aleatoare;
    \item \texttt{Sorted} -- valori deja ordonate crescător;
    \item \texttt{Reverse} -- valori ordonate descrescător;
    \item \texttt{NearlySorted} -- vector aproape sortat;
    \item \texttt{FewUnique} -- vector cu puține valori distincte.
\end{itemize}

Dimensiunile utilizate în benchmark-ul din proiect sunt:

\begin{center}
\texttt{10, 20, 50, 100, 500, 1000, 5000, 10000}
\end{center}

Pentru dimensiunile mici au fost efectuate multe iterații, pentru a obține o medie relevantă, iar pentru dimensiunile mai mari numărul de iterații a fost redus. În plus, pentru fiecare execuție se verifică dacă vectorul rezultat este într-adevăr sortat corect.

\section{Analiza rezultatelor}

Rezultatele experimentale confirmă, în linii mari, teoria complexității algoritmilor. Algoritmii pătratici devin rapid ineficienți pe măsură ce dimensiunea datelor crește, în timp ce algoritmii moderni își păstrează performanța într-o zonă mult mai bună.

Bubble Sort și Selection Sort sunt utile în principal pentru scop didactic. Chiar dacă implementările sunt corecte și simple, timpii de execuție cresc foarte repede pentru volume mai mari de date. Insertion Sort rămâne totuși interesant, deoarece pe date aproape sortate se comportă semnificativ mai bine decât în scenariile aleatoare.

Merge Sort și Quick Sort oferă rezultate solide pentru dimensiuni mari. Merge Sort este apreciat pentru predictibilitate și stabilitate, în timp ce Quick Sort are de obicei timpi mai buni în practică. Heap Sort rămâne o variantă robustă, mai ales atunci când se dorește o metodă cu memorie suplimentară redusă.

Cele trei metode adăugate ulterior proiectului completează foarte bine analiza:

\begin{itemize}
    \item Shell Sort este o soluție intermediară foarte utilă între algoritmii simpli și cei moderni;
    \item Radix Sort este extrem de competitiv pentru date numerice întregi;
    \item IntroSort oferă un reper practic foarte valoros, deoarece este apropiat de ceea ce se folosește în programele reale.
\end{itemize}

În ansamblu, se observă că nu există un algoritm universal perfect pentru orice situație. Alegerea depinde de dimensiunea datelor, de distribuția acestora și de constrângerile legate de memorie.

\section{Concluzii}

Acest proiect demonstrează clar diferențele dintre algoritmii de sortare clasici și cei performanți. Algoritmii simpli precum Bubble Sort, Selection Sort și Insertion Sort sunt potriviți pentru învățare și pentru exemple de dimensiuni mici, însă pentru volume mari devin ineficienți.

Merge Sort, Quick Sort, Heap Sort, Shell Sort, Radix Sort și IntroSort oferă rezultate semnificativ mai bune în practică. Dintre acestea, IntroSort și Quick Sort sunt foarte potrivite pentru uz general, Merge Sort este foarte stabil, iar Radix Sort este foarte bun atunci când se lucrează exclusiv cu valori întregi.

Din perspectivă academică, proiectul este valoros deoarece îmbină partea teoretică a analizei algoritmilor cu o validare experimentală concretă. Lucrarea poate fi folosită atât ca material de studiu, cât și ca bază pentru dezvoltări viitoare, precum adăugarea unor noi algoritmi, generarea automată de grafice sau extinderea testelor pentru seturi mai mari de date.

\section{Bibliografie}

\begin{thebibliography}{9}

\bibitem{cormen}
Thomas H. Cormen, Charles E. Leiserson, Ronald L. Rivest, Clifford Stein,
\textit{Introduction to Algorithms},
The MIT Press, ediția a 4-a, 2022.

\bibitem{sedgewick}
Robert Sedgewick, Kevin Wayne,
\textit{Algorithms},
Addison-Wesley Professional, ediția a 4-a, 2011.

\bibitem{knuth}
Donald E. Knuth,
\textit{The Art of Computer Programming, Volume 3: Sorting and Searching},
Addison-Wesley, ediția a 2-a, 1998.

\bibitem{shell}
D. L. Shell,
\textit{A High-Speed Sorting Procedure},
Communications of the ACM, vol. 2, nr. 7, 1959, pp. 30--32.

\bibitem{musser}
David R. Musser,
\textit{Introspective Sorting and Selection Algorithms},
Software: Practice and Experience, vol. 27, nr. 8, 1997, pp. 983--993.

\bibitem{bentley}
Jon L. Bentley, M. Douglas McIlroy,
\textit{Engineering a Sort Function},
Software: Practice and Experience, vol. 23, nr. 11, 1993, pp. 1249--1265.

\end{thebibliography}

\end{document}
